\section{Anwendungen und Chancen intelligenter Assistenten}

\subsection{Anwendungsmöglichkeiten im privaten und beruflichen Alltag}

Intelligente Assistenten haben sich im privaten und beruflichen Alltag zu vielseitigen Schnittstellen zwischen Nutzenden und digitalen Diensten entwickelt. Bereits etablierte Systeme wie Siri, Alexa oder Google Assistant ermöglichen unter anderem die Informationssuche, die Verwaltung von Terminen und Erinnerungen sowie die sprachgesteuerte Kommunikation \cite{silva2020intelligent}. Im privaten Umfeld kommt insbesondere die Verknüpfung mit Smart-Home-Systemen hinzu. Sprachassistenten können dabei als zentrale Benutzerschnittstelle zu vernetzten Geräten und Diensten dienen und beispielsweise die Steuerung von Haushaltsgeräten oder die Organisation alltäglicher Routinen unterstützen \cite{minder2023voice}. Im beruflichen Kontext lassen sich vergleichbare Funktionen für die Terminplanung, Kommunikation und Informationsbeschaffung einsetzen, wodurch wiederkehrende organisatorische Tätigkeiten zunehmend an digitale Assistenzsysteme delegiert werden können \cite{silva2020intelligent}.

Gleichzeitig zeichnet sich ein grundlegender Wandel ihres Funktionsumfangs ab. Während klassische Sprachassistenten überwiegend einzelne, explizit formulierte Befehle ausführen, ermöglichen aktuelle Entwicklungen rund um Large Language Models zunehmend Systeme, die Anweisungen interpretieren, mehrstufige Aufgaben planen und externe Werkzeuge in die Bearbeitung einbeziehen können \cite{chowa2026language}. Damit entwickelt sich das Konzept des intelligenten Assistenten perspektivisch von einer sprachgesteuerten Benutzerschnittstelle hin zu einer Assistenz, die komplexere Aufgaben über verschiedene Anwendungen und Dienste hinweg koordinieren kann.